% IEEE conference paper for the AFML course project.
% Overleaf: upload main.tex plus a figures/ folder holding a copy of results/figures/*.png, and compile
% with pdfLaTeX. Compiled inside the repository, the figures are found in ../../results/figures/.
% Every number comes from a file in results/; results/README.md says which. Nothing here is estimated.
\documentclass[conference]{IEEEtran}
\usepackage{amsmath,amssymb}
\usepackage{graphicx}
\usepackage{booktabs}
\usepackage{xcolor}
\usepackage{url}
\usepackage[hidelinks]{hyperref}

\newcommand{\todo}[1]{\textcolor{red}{[TODO: #1]}}
% Includes figures/<file> or ../../results/figures/<file>, otherwise shows a red placeholder box.
\newcommand{\figfile}[2][width=\linewidth]{%
  \IfFileExists{figures/#2}{\includegraphics[#1]{figures/#2}}{%
  \IfFileExists{../../results/figures/#2}{\includegraphics[#1]{../../results/figures/#2}}{%
    \fbox{\parbox{0.92\linewidth}{\centering\vspace{1.5em}%
      \textcolor{red}{figure pending: \texttt{\detokenize{#2}}}\vspace{1.5em}}}}}}

\begin{document}

\title{Swapping Attention for Selective State Spaces:\\
Training and Evaluating a Mamba-Style Small Language Model}

\author{\IEEEauthorblockN{Krishna G Pai, Dikshith Varma M N, Gagandeep N, Mallikarjuna Rao R V}
\IEEEauthorblockA{Department of Computer Science and Engineering (AI \& ML), PES University, Bengaluru, India\\
PES1UG24AM372, PES1UG24AM363, PES1UG24AM367, PES1UG24AM155}}

\maketitle

\begin{abstract}
State-space models (SSMs) such as Mamba replace attention with a recurrence whose state has a fixed
size, so the cost of generating a token does not grow with the context. We compare a Transformer and a
Mamba-style selective SSM as small language models under matched conditions: about 47M parameters
each, the same FineWeb-Edu data and tokenizer, and the same budget of 340M training tokens, with three
seeds per architecture. A third model, Mamba-LTI, removes the input-dependent (selective) part of Mamba
and nothing else. We measure test loss, associative recall (the synthetic MQAR task, and natural-text
loss split into recall-hit and other tokens), loss beyond the 1,024-token training length on PG-19
books, and inference latency and memory with cached decoding. Finally, we instruction-tune both models
on the same instruction data and compare their answers before and after.
Mamba reaches a test loss of 3.484 against 3.531 for the Transformer (perplexity 32.6 against 34.1), and
every Mamba seed beats every Transformer seed. Removing selectivity costs 0.176 nats. The Transformer is
better on the 16.8\% of test tokens whose answer can be copied from the context (1.39 against 1.76);
Mamba is better on all other tokens, and its loss on PG-19 books rises only from 3.40 to 3.88 between
the training length and 16K tokens, where the Transformer's rises to 5.62. With the official
\texttt{mamba-ssm} kernel, Mamba's decoding memory stays at 116 MiB at batch 8 and 32K tokens, where
the Transformer's reaches 6,260 MiB; but on a T4 the same Mamba run takes twice as long to train, so
at equal training time the Transformer is ahead. After instruction tuning, both models answer simple
questions and their response loss differs by 0.04 nats. All training ran on free Kaggle T4 GPUs; the
code, configurations and per-run logs are public.
\end{abstract}

\begin{IEEEkeywords}
state-space models, Mamba, Transformer, small language models, selectivity, associative recall,
instruction tuning
\end{IEEEkeywords}

% ---------------------------------------------------------------------------------------------------
\section{Introduction}

A Transformer~\cite{vaswani2017attention} reads its context through attention: each new token is
compared with every earlier token. Processing a prompt of $L$ tokens therefore costs $O(L^2)$ per
layer, and the key-value (KV) cache that generation needs grows linearly with $L$. State-space models
process the sequence with a linear recurrence instead. Their state has a fixed size, so each generated
token costs the same time and memory however long the context is.

Early SSMs such as S4~\cite{gu2022s4} are \emph{linear time-invariant} (LTI): every token is processed
with the same dynamics. This allows fast training as a convolution, but the model cannot decide from
the content what to keep and what to ignore. Mamba~\cite{gu2023mamba} makes the recurrence
\emph{selective}: its step size and its input and output maps are computed from the current token. Its
authors report quality equal to or better than Transformers of the same size, in experiments that
start at 130M parameters and use far larger training budgets than a student project can afford.
Zoology~\cite{arora2023zoology} found that much of the remaining gap between attention and
attention-free models comes from associative recall: retrieving a value seen earlier in the context
when its key appears again.

\textbf{Research question.} At the same parameter count, the same data and the same token budget, how
does a Mamba-style small language model compare with a Transformer in quality and in inference
efficiency, and what does selectivity contribute?

\textbf{Contributions.}
\begin{enumerate}
  \item A matched comparison of a Transformer and a Mamba model with about 47M parameters each, three
        seeds per architecture, 340M training tokens of FineWeb-Edu, trained entirely on free T4 GPUs.
  \item A selectivity ablation (Mamba-LTI) with the same size, data and training.
  \item A recall analysis on synthetic data (MQAR) and on natural text (recall-hit tokens).
  \item Loss beyond the training length, up to 16,384 tokens of PG-19 books.
  \item Inference latency and memory with a KV cache against a recurrent state, reported separately for
        the official Mamba kernel and for our pure-PyTorch scan.
  \item Instruction fine-tuning of both models with identical data and settings, with response loss and
        answer quality before and after.
  \item Public code, including a memory-bounded chunked selective scan in pure PyTorch.
\end{enumerate}

% ---------------------------------------------------------------------------------------------------
\section{Background}

\subsection{Attention}
Causal self-attention computes $\mathrm{softmax}(QK^\top/\sqrt{d})\,V$ with future positions masked.
Prefill over $L$ tokens costs $O(L^2 d)$ per layer. Decoding one token with a KV cache costs $O(Ld)$
per layer and stores $2Ld$ values per layer.

\subsection{State-space layers}
A diagonal SSM channel keeps a state $h_t$ of size $N$ (here $N = 16$) and updates it per token:
\begin{align}
  h_t &= \exp(\Delta A) \odot h_{t-1} + \Delta B\, u_t, \label{eq:ssm}\\
  y_t &= C\, h_t + D\, u_t .
\end{align}
$A$ is diagonal and negative, so $\exp(\Delta A)$ lies between 0 and 1 and acts as a forget factor. In
an LTI model, $\Delta$, $B$ and $C$ are the same at every position.

\subsection{Selectivity}
Mamba computes $\Delta_t = \mathrm{softplus}(W_\Delta x_t + b)$, $B_t = W_B x_t$ and $C_t = W_C x_t$ from
the current token. A large $\Delta_t$ drives $\exp(\Delta_t A)$ towards 0, so the state is overwritten
with the current token. A small $\Delta_t$ keeps the state and ignores the token. A channel's memory
horizon is therefore about $\tau = 1/(\Delta_t |A|)$ tokens. Because the dynamics change from token to
token, training cannot use a convolution and needs a scan over the sequence. Mamba-2~\cite{dao2024mamba2}
restricts $A$ further so that the scan becomes matrix multiplications; we use the Mamba-1 formulation.

\subsection{Associative recall}
In multi-query associative recall (MQAR)~\cite{arora2023zoology}, a sequence lists key-value pairs and
later repeats some keys; the model must output the matching values. Attention can look a key up
directly. A recurrent model must already have stored the pair in its fixed-size state.

\subsection{Instruction fine-tuning}
A pretrained language model only continues text. Instruction fine-tuning~\cite{wei2022flan} trains it
further on pairs of an instruction and a good response, so that it answers a new instruction instead of
continuing it. The loss is usually taken on the response tokens only: the instruction is context, not
text the model should learn to produce.

% ---------------------------------------------------------------------------------------------------
\section{Models}

\subsection{Transformer}
A pre-norm decoder with RMSNorm~\cite{zhang2019rmsnorm}, rotary position embeddings
(RoPE, base 10,000)~\cite{su2021roformer}, a SwiGLU MLP~\cite{shazeer2020glu} with hidden size $8/3$ of
the width rounded up to a multiple of 64, PyTorch scaled-dot-product attention with 8 heads of size 64,
and tied input and output embeddings. Generation uses a KV cache.

\subsection{Mamba}
Each layer is RMSNorm, a Mamba-1 block and a residual connection. The block projects its input to a
signal $x$ and a gate $z$ of width $d_\text{inner}$, applies a causal depthwise convolution of width 4
and SiLU to $x$, runs the selective SSM (state size 16, $\Delta$ projection of rank 32), gates the
result with $\mathrm{SiLU}(z)$ and projects back to the model width. $A$ is initialised to
$-1,\dots,-16$ (S4D-real) and the $\Delta$ bias so that $\mathrm{softplus}$ of it lies between 0.001
and 0.1. A Mamba layer has no MLP, so the Mamba models have two layers for each Transformer layer, as
in~\cite{gu2023mamba}. Generation uses the recurrent state (convolution state and SSM state), whose
size does not depend on the context length.

\subsection{Mamba-LTI ablation}
$\Delta$, $B$ and $C$ are replaced by learned constants that do not depend on the input. Everything
else is identical, so the ablation removes selectivity and nothing more. Without the three input
projections the block has fewer parameters, so its $d_\text{inner}$ is matched separately.

\subsection{Parameter matching}
All models share the width (512), the vocabulary (16,384) and the tied embedding. The Mamba
$d_\text{inner}$ is found by binary search so that non-embedding parameters are within 0.5\% of the
Transformer's (Table~\ref{tab:params}); the gaps are 0.22\% for Mamba and 0.20\% for Mamba-LTI. The
counts in the table are read from the trained runs' summaries.

\begin{table}[t]
\centering
\caption{Model sizes (scale S, width 512). Counts from the trained runs' summaries.}
\label{tab:params}
\begin{tabular}{lrrrr}
\toprule
Model & Layers & $d_\text{inner}$ & Non-embedding & Total \\
\midrule
Transformer & 12 & -- & 38,547,968 & 46,936,576 \\
Mamba & 24 & 968 & 38,461,760 & 46,850,368 \\
Mamba-LTI & 24 & 1,032 & 38,626,880 & 47,015,488 \\
\bottomrule
\end{tabular}
\end{table}

\subsection{Scan implementations}
We implement four interchangeable scans: a step-by-step reference, a Hillis--Steele parallel scan over
the whole sequence, a chunked scan, and the fused CUDA kernel of the official \texttt{mamba-ssm}
package. The parallel scan keeps tensors of shape (batch, length, $d_\text{inner}$, $N$) for
backpropagation and ran out of memory for a full-size model at 1,024 tokens in our tests. The chunked
scan runs the parallel scan inside chunks of 64 tokens with activation checkpointing and carries the
state between chunks, so it stores only the block inputs and one state per chunk. Unit tests check that
every backend gives the same outputs and gradients as the reference. The scan always runs in fp32.
Every speed number below names the backend (\texttt{cuda} or \texttt{chunked}), because a slow
pure-PyTorch scan is a property of our implementation, not of the architecture. All Mamba and
Mamba-LTI training runs and all fine-tuning runs used the \texttt{cuda} kernel, as recorded in each
run's summary. The \texttt{chunked} scan was used on the GPU only in the smoke test and in the
efficiency benchmark, where it is timed next to the kernel.

% ---------------------------------------------------------------------------------------------------
\section{Data and Training}

\subsection{Data}
The training text is FineWeb-Edu~\cite{penedo2024fineweb}, subset \texttt{sample-10BT}, pinned to one
revision of the dataset repository. We stream the subset through a 10,000-document shuffle buffer
(seed 1234), drop exact duplicates by the SHA-256 hash of the stripped text (12,502 documents), and
assign each document to the test, validation or training split from its hash (1\%, 1\% and 98\%), so
the split does not depend on order. A byte-level BPE tokenizer with 16,384 tokens is fitted on 100,000
training documents only, and documents are separated by an end-of-document token. The build stops at
2B training tokens (Table~\ref{tab:data}). Each run draws its training tokens as random 1,024-token
windows from the training split.

For long-context evaluation we use the first 50 PG-19~\cite{rae2020compressive} test books that have
at least 16,385 tokens under our tokenizer, read from the Parquet copy \texttt{emozilla/pg19} of the
official release.

For instruction fine-tuning we use Alpaca-cleaned (\texttt{yahma/alpaca-cleaned}, one pinned revision,
CC BY 4.0), a cleaned version of the Stanford Alpaca pairs~\cite{taori2023alpaca}, which were generated
with an OpenAI model following Self-Instruct~\cite{wang2023selfinstruct}. We drop exact duplicates and
split the pairs by the hash of their text into test (2\%), validation (2\%) and training. Each example
is written as \texttt{\#\#\# Instruction:}, the instruction, an optional \texttt{\#\#\# Input:} block,
\texttt{\#\#\# Response:}, the response and the end-of-document token, and is tokenized with the
pretraining tokenizer; examples longer than 1,025 tokens are dropped. Of the 51,760 pairs, 4 are
duplicates and 30 are too long. The splits hold 49,570 training, 1,081 validation and 1,075 test
examples, on average 197 tokens each; the test split has 170,888 response tokens.

\begin{table}[t]
\centering
\caption{FineWeb-Edu splits after deduplication.}
\label{tab:data}
\begin{tabular}{lrr}
\toprule
Split & Documents & Tokens \\
\midrule
Train & 1,809,687 & 2,000,000,330 \\
Validation & 8,945 & 10,000,935 \\
Test & 9,159 & 10,000,062 \\
\bottomrule
\end{tabular}
\end{table}

\subsection{Training setup}
Table~\ref{tab:train} lists the settings, which are identical for all models except the learning rate
and the micro-batch split. If a micro-batch does not fit in GPU memory, the trainer halves it and
doubles gradient accumulation, so the tokens per step never change.

\begin{table}[t]
\centering
\caption{Training settings.}
\label{tab:train}
\begin{tabular}{@{}p{0.34\linewidth}p{0.58\linewidth}@{}}
\toprule
Optimizer & AdamW~\cite{loshchilov2019adamw}, $\beta = (0.9, 0.95)$, weight decay 0.1 on weight matrices only (none on norms, biases, embeddings, $A$, $D$ or the $\Delta$ bias) \\
Schedule & linear warmup for 51 steps (2\%), cosine decay to 10\% of the peak \\
Gradient clipping & 1.0 \\
Context & 1,024 tokens \\
Tokens per step & 131,072 (Transformer $16 \times 8$, Mamba $8 \times 16$ micro-batch $\times$ accumulation) \\
Training tokens & 339,869,696 (2,593 steps), about 7.2 per parameter \\
Precision & fp16 with loss scaling (the T4 has no native bf16); scan in fp32 \\
Validation & 2M tokens at step 0, every 300 steps and at the last step; the best checkpoint is kept \\
Seeds & Transformer and Mamba: 0, 1, 2; Mamba-LTI: 0, 1 \\
\bottomrule
\end{tabular}
\end{table}

\subsection{Token budget}
We first planned about 20 tokens per parameter, following Chinchilla~\cite{hoffmann2022chinchilla}. A
short GPU smoke test on the free Kaggle T4 measured 10,370 tokens/s for the Transformer and
10,815 tokens/s for Mamba with the \texttt{cuda} kernel, which put a full run at about 25 hours. Mamba
with our \texttt{chunked} scan reached only 380 tokens/s, about 28 times slower than the kernel. We
therefore cut the budget to 340M tokens so that the slower model finishes within one 11-hour free
session, and reduced the step size so that the runs still take about 2,600 optimizer steps. Over the
full runs the Transformer was much faster than the smoke test suggested: seed 0 averaged
23,221 tokens/s (4.1 hours), against 11,198 tokens/s (8.4 hours) for Mamba with the \texttt{cuda}
kernel.

\subsection{Learning rate}
Each architecture was trained for 30M tokens (228 steps) with the main runs' step size and their
51 warmup steps, at learning rates from $1.5\times10^{-4}$ to $2.4\times10^{-3}$ (Table~\ref{tab:lr}).
The rate with the lowest validation loss is used for the full runs, with one tie rule: rates within
0.01 of the lowest loss count as tied and the lowest tied rate is used, because the full runs are about
11 times longer, the best rate tends to fall as training gets longer, and a lower rate is less likely
to diverge in fp16. We added this rule after seeing the sweep, in which two Mamba rates came out almost
equal; we state it so the choice is transparent. The Transformer's best rate, $6\times10^{-4}$, lies
inside the grid. For Mamba, $2.4\times10^{-3}$ gave the lowest loss, only 0.0008 below
$1.2\times10^{-3}$, so the tie rule chose $1.2\times10^{-3}$. Because Mamba's lowest loss sits at the
top of the grid, a higher rate might have done slightly better in these short runs; we did not test
one. Mamba-LTI uses the Mamba rate.

A first sweep used 2\% of each sweep run's own length as warmup, only 4 steps. Two Mamba rates diverged
and the Transformer's best rate was the lowest one tried. Because the main runs warm up for 51 steps,
that sweep was biased against high rates, so it was discarded and repeated with the main-run warmup.
Its log is kept in the repository.

\begin{table}[t]
\centering
\caption{Validation loss after 30M tokens in the learning-rate sweep. No run diverged. Bold: the rate used.}
\label{tab:lr}
\begin{tabular}{lrr}
\toprule
Learning rate & Transformer & Mamba \\
\midrule
$1.5\times10^{-4}$ & 6.022 & -- \\
$3\times10^{-4}$ & 5.644 & 5.602 \\
$6\times10^{-4}$ & \textbf{5.413} & 5.201 \\
$1.2\times10^{-3}$ & 5.619 & \textbf{4.924} \\
$2.4\times10^{-3}$ & 5.744 & 4.923 \\
\bottomrule
\end{tabular}
\end{table}

\subsection{Hardware}
All runs used free Kaggle notebooks with two NVIDIA T4 GPUs (16 GB each), running two independent runs
at a time, one per GPU. The eight main runs were split over the four team members' accounts. Training
saves a full checkpoint (model, optimizer, loss scaler and random states) before the 12-hour session
limit, so a run can resume in a later session.

\subsection{Instruction fine-tuning}
The seed-0 Transformer and Mamba are fine-tuned from their best pretraining checkpoints with identical
settings. Every weight is trained. The loss is cross-entropy on the response tokens and the closing
end-of-document token only. We use AdamW with $\beta = (0.9, 0.95)$ and no weight decay, 64 examples per
step (micro-batch 8, accumulation 8, right-padded and grouped by length), 2 epochs, linear warmup over
3\% of the steps and cosine decay to 10\% of the peak, gradient clipping at 1.0, and fp16 with the scan
in fp32. Each model is fine-tuned at learning rates $3\times10^{-5}$, $10^{-4}$, $3\times10^{-4}$ and
$10^{-3}$ (the last was added because both models were best at $3\times10^{-4}$, the edge of the first
grid); the rate is chosen on validation response loss with the same tie rule as in pretraining, and the checkpoint
with the lowest validation loss is kept. Both models were best at $3\times10^{-4}$. For Mamba,
$10^{-3}$ came within 0.01 of it (2.195 against 2.185 validation response loss), so the tie rule again
picked the lower rate. Each fine-tuning run takes 1,550 steps: 14 to 15 minutes for the Transformer and
39 to 41 minutes for Mamba with the \texttt{cuda} kernel, on one T4.

% ---------------------------------------------------------------------------------------------------
\section{Evaluation Protocol}

The test split is scored once, with the checkpoint chosen on validation; nothing is selected on it.

\textbf{E1, quality.} Test loss (nats per token) and perplexity on 5M test tokens in 1,024-token
windows; mean and standard deviation over seeds.

\textbf{E2, selectivity.} Mamba against Mamba-LTI on E1, and on selective copying~\cite{gu2023mamba}:
16 data tokens placed at random among noise tokens must be output in order; width 64, sequence lengths
64 and 256, learning rates $10^{-3}$ and $3\times10^{-3}$.

\textbf{E3, associative recall.} MQAR at length 128 with 16 key-value pairs, widths 64, 128 and 256,
learning rates $3\times10^{-4}$, $10^{-3}$, $3\times10^{-3}$ and $10^{-2}$. Both synthetic tasks use
2 Transformer layers against 4 Mamba layers, 40,000 steps at batch 64, one seed, accuracy on 1,000
held-out sequences every 1,000 steps, and stop early at 99.5\%. We report every learning rate, so no
model is favoured. A first round of 4,000-step runs over lengths 64 to 512 is kept in the repository;
it was too short, and its conclusions differ from the long runs (Section~\ref{sec:recall}). On natural text,
a test token is a \emph{recall hit} if the pair (previous token, token) already occurred earlier in
the same window~\cite{arora2023zoology}; we report the loss on hits and on all other tokens separately.

\textbf{E4, length extrapolation.} The first 16,384 tokens of each PG-19 book in one forward pass, with
per-token loss averaged in position bins that double in size from 128 tokens, for the seed-0 models.
RoPE is applied beyond 1,024 positions without rescaling.

\textbf{E5, efficiency.} Randomly initialised models with the trained configurations (speed does not
depend on the weights), fp16 on one T4: prefill latency, throughput and peak memory, and per-token
decode latency and peak memory with a KV cache (Transformer) or the recurrent state (Mamba), for 512 to
32,768 tokens at batch 1 and 8. Prefill is measured as generation runs it: the prompt fills the cache or
state and only the last position is scored, and its peak memory includes the weights and the cache. Timing uses CUDA events, the median of 10 runs after 3 warm-up runs. Mamba is timed with
both the \texttt{cuda} and the \texttt{chunked} backend. Out-of-memory points are missing, and once one
prefill takes over 20 s, longer lengths at that batch size are skipped.

\textbf{E6, mechanism.} $\Delta_t$ for every layer and token of the seed-0 Mamba on the first 256 test
tokens, the memory horizons it implies, and the seed-0 Transformer's attention maps on the same text.

\textbf{E7, instruction following.} For each model before and after fine-tuning: the loss on the
response tokens of every test example (prompt tokens are not scored); greedy answers of at most 128 new
tokens to the first 100 test prompts, scored by ROUGE-L F1~\cite{lin2004rouge} against the reference
answer and by the share of answers that end with the end-of-document token; and answers to eight
hand-written prompts, shown as examples. One seed.

% ---------------------------------------------------------------------------------------------------
\section{Results}

\subsection{Language-modelling quality (E1)}
Table~\ref{tab:main} gives the test results and Fig.~\ref{fig:curves} the validation curves. Mamba ends
at a test loss of 3.484, the Transformer at 3.531, a gap of 0.046 nats (perplexity 32.6 against 34.1).
The standard deviation over seeds is at most 0.006 for every model, and the worst Mamba seed (3.490) is
below the best Transformer seed (3.525), so the gap is not seed noise. Mamba is ahead at every
validation point after step 0: after 300 steps (39M tokens) it is at 4.53 to 4.56 against 4.89 to 4.93,
and the gap narrows to about 0.04 by the end (final validation loss 3.509 against 3.550, means over
seeds). Mamba-LTI ends at 3.661, 0.176 above Mamba and 0.130 above the
Transformer.

\begin{table}[t]
\centering
\caption{Test results on 5M FineWeb-Edu test tokens, mean $\pm$ standard deviation over seeds. Hit:
recall-hit tokens (16.8\% of the test tokens); other: all remaining tokens. Perplexity is
$\exp$ of the mean loss.}
\label{tab:main}
\begin{tabular}{@{}lrrrrr@{}}
\toprule
Model & Seeds & Test loss & Ppl & Hit loss & Other loss \\
\midrule
Transformer & 3 & $3.531 \pm 0.005$ & 34.14 & \textbf{1.387} & 3.964 \\
Mamba & 3 & $\mathbf{3.484 \pm 0.005}$ & \textbf{32.59} & 1.762 & \textbf{3.832} \\
Mamba-LTI & 2 & $3.661 \pm 0.006$ & 38.88 & 2.416 & 3.912 \\
\bottomrule
\end{tabular}
\end{table}

\begin{figure}[t]
\centering
\figfile{fig1_training_curves.png}
\caption{Validation loss against training tokens and training FLOPs. The band is the range over seeds.}
\label{fig:curves}
\end{figure}

\textbf{Same time instead of same tokens.} The token budget is equal, but the training time is not
(Fig.~\ref{fig:time}). The Transformer runs finished in 3.98 to 4.44 hours (mean 4.16 h) at a mean
final validation loss of 3.550. Mamba, with the \texttt{cuda} kernel, needed 8.15 to 8.50 hours. At the
4.16-hour mark the Mamba runs had seen 157M tokens and stood at 3.770 to 3.792; they reached 3.550 after
6.70 to 7.25 hours. By our FLOP count Mamba does slightly less arithmetic per token
($9.9\times10^{16}$ against $1.09\times10^{17}$ FLOPs for the whole run), so the time gap comes from
how well the two implementations use the T4, not from the amount of work: the Transformer runs at
21,000 to 24,000 tokens/s against 11,000 to 12,000 for Mamba, and Mamba's micro-batch had to be halved
to 8 to fit in memory. Two caveats. A run's loss partway through is worse than what a run planned for
that shorter budget would reach, because its learning rate has not yet decayed, so the 4-hour Mamba
value is pessimistic. And the result is specific to this GPU, fp16 and this kernel. Still, for a
student with a fixed number of free GPU hours, the Transformer gave the lower loss per hour, and Mamba
the lower loss per token.

\begin{figure}[t]
\centering
\figfile{fig9_time_budget.png}
\caption{Validation loss against training time on one T4 (fp16; Mamba with the \texttt{mamba-ssm}
kernel). Every run sees 340M tokens; the dot marks the end of each run.}
\label{fig:time}
\end{figure}

\subsection{Selectivity (E2)}
Replacing the input-dependent $\Delta$, $B$ and $C$ by learned constants raises the test loss from 3.484
to 3.661 (Table~\ref{tab:main}). The split in the last two columns shows where the loss goes: on
recall-hit tokens Mamba-LTI is 0.65 nats worse than Mamba, on all other tokens only 0.08. On the other
tokens Mamba-LTI (3.912) is even slightly better than the Transformer (3.964); almost its whole deficit
is on tokens that must be copied from the context.

Selective copying makes the same point directly (Fig.~\ref{fig:copy}, Table~\ref{tab:copy}). Mamba
reaches 99\% at both lengths and both learning rates, within 9,000 to 26,000 steps. Mamba-LTI reaches
99\% only at $3\times10^{-3}$, after 28,000 to 29,000 steps; at $10^{-3}$ it stops at 54.2\% (length 64)
and 12.9\% (length 256) after 40,000 steps. The Transformer solves length 64 but stops at 91.9\% to
93.3\% at length 256. So an LTI model can learn the task, but in our runs it needed a higher learning
rate and two to three times as many steps as the selective model.

\begin{table}[t]
\centering
\caption{Selective copying, best held-out accuracy in 40,000 steps (one seed), and the first step at
which accuracy reached 99\%.}
\label{tab:copy}
\begin{tabular}{@{}llrrrr@{}}
\toprule
& & \multicolumn{2}{c}{Length 64} & \multicolumn{2}{c}{Length 256} \\
Model & LR & Acc. & Step & Acc. & Step \\
\midrule
Transformer & $10^{-3}$ & 99.6\% & 31k & 91.9\% & -- \\
Transformer & $3\times10^{-3}$ & 99.7\% & 10k & 93.3\% & -- \\
Mamba & $10^{-3}$ & 99.7\% & 15k & 99.6\% & 26k \\
Mamba & $3\times10^{-3}$ & 99.7\% & 9k & 99.6\% & 13k \\
Mamba-LTI & $10^{-3}$ & 54.2\% & -- & 12.9\% & -- \\
Mamba-LTI & $3\times10^{-3}$ & 99.5\% & 29k & 99.7\% & 28k \\
\bottomrule
\end{tabular}
\end{table}

\begin{figure}[t]
\centering
\figfile{fig2_selective_copy.png}
\caption{Selective copying, 40,000 training steps, one seed: held-out accuracy at two learning rates
(line style). Runs stop at 99.5\%.}
\label{fig:copy}
\end{figure}

\subsection{Associative recall (E3)}
\label{sec:recall}
On natural text the picture reverses (Fig.~\ref{fig:recall}). On the 841,199 recall-hit tokens the
Transformer's loss is 1.387 against Mamba's 1.762, a gap of 0.375 nats in the Transformer's favour,
consistent across seeds (1.36 to 1.42 against 1.75 to 1.77). On the other 4,157,969 tokens Mamba leads
by 0.13 (3.832 against 3.964). The hit tokens are 16.8\% of the test set, so the two parts combine to
Mamba's overall lead of 0.046. In other words, the Transformer's advantage is concentrated on tokens it
can look up, and Mamba's advantage is spread over everything else.

On MQAR (Fig.~\ref{fig:mqar}, Table~\ref{tab:mqar}), no model solved width 64 within 40,000 steps; the
best is Mamba at 19.6\%. At widths 128 and 256 both the Transformer and Mamba reached 99.5\% or more,
each at exactly one of the four learning rates: the Transformer at $10^{-3}$ (width 128, step 24,000)
and $10^{-2}$ (width 256, step 17,000), Mamba at $10^{-2}$ (width 128, step 29,000) and
$3\times10^{-3}$ (width 256, step 18,000). Mamba-LTI never passed 20.6\%. Of the 32 runs that failed, 20 peak between 18\% and
22\%. Each key is queried once and its value follows the query, so a model that guesses
uniformly among the values not yet asked would score $(1 + 1/2 + \dots + 1/16)/16 = 0.211$; we read the
plateau that way but did not test it. At this size and training length we therefore cannot say that
either architecture is better at MQAR; the result is that both can solve it, LTI cannot, and that the
outcome depends strongly on the learning rate.

The earlier 4,000-step runs told a different story. On selective copying the Transformer reached
60\% to 94\% while both Mamba models stayed at 6\% to 13\%, and on MQAR at length 128 the Transformer
reached 21\% against at most 9\% for Mamba. With ten times more steps, Mamba caught up on MQAR and
overtook the Transformer on selective copying at length 256. We mention this because a short synthetic
run can point the wrong way.

\begin{table}[t]
\centering
\caption{MQAR at length 128 with 16 pairs: best held-out accuracy in 40,000 steps (one seed). Bold:
solved (at least 99.5\%).}
\label{tab:mqar}
\begin{tabular}{@{}llrrr@{}}
\toprule
Model & LR & Width 64 & Width 128 & Width 256 \\
\midrule
Transformer & $3\times10^{-4}$ & 17.8\% & 20.6\% & 21.9\% \\
Transformer & $10^{-3}$ & 18.1\% & \textbf{99.5\%} & 21.9\% \\
Transformer & $3\times10^{-3}$ & 3.1\% & 5.7\% & 21.8\% \\
Transformer & $10^{-2}$ & 18.1\% & 5.4\% & \textbf{99.6\%} \\
\midrule
Mamba & $3\times10^{-4}$ & 9.9\% & 19.7\% & 21.1\% \\
Mamba & $10^{-3}$ & 18.6\% & 19.9\% & 21.1\% \\
Mamba & $3\times10^{-3}$ & 19.6\% & 21.3\% & \textbf{99.5\%} \\
Mamba & $10^{-2}$ & 16.0\% & \textbf{99.6\%} & 20.2\% \\
\midrule
Mamba-LTI & $3\times10^{-4}$ & 0.2\% & 2.9\% & 12.3\% \\
Mamba-LTI & $10^{-3}$ & 15.5\% & 15.3\% & 18.3\% \\
Mamba-LTI & $3\times10^{-3}$ & 19.2\% & 16.6\% & 20.6\% \\
Mamba-LTI & $10^{-2}$ & 18.2\% & 16.9\% & 20.4\% \\
\bottomrule
\end{tabular}
\end{table}

\begin{figure*}[t]
\centering
\figfile[width=\textwidth]{fig3_mqar.png}
\caption{MQAR at length 128 with 16 key-value pairs: accuracy during training for the best of four
learning rates per model and width, one seed. The dotted line is 0.21.}
\label{fig:mqar}
\end{figure*}

\begin{figure}[t]
\centering
\figfile{fig4_recall_split.png}
\caption{Test loss on recall-hit tokens and on all other tokens.}
\label{fig:recall}
\end{figure}

\subsection{Length extrapolation (E4)}
Fig.~\ref{fig:long} shows the loss by position on the 50 PG-19 books, for the seed-0 models. Inside the
training length the two main models are close, and the Transformer is slightly better: 3.291 against
3.351 at positions 256 to 512, and 3.396 against 3.399 at 512 to 1,024. Past 1,024 tokens, where RoPE
sees positions it was never trained on, the Transformer's loss jumps to 4.237 (1K to 2K), 5.260 (2K to
4K) and 5.619 (8K to 16K). Mamba's loss also rises, to 3.589, 3.809 and 3.879 in the same bins, so it
does not keep improving with more context either, but it degrades by 0.48 nats where the Transformer
degrades by 2.22. Mamba-LTI behaves like Mamba (3.702 at 512 to 1,024, 4.102 at 8K to 16K), which is
expected: a recurrence has no position index to go out of range. We did not try a RoPE scaling method,
which might reduce the Transformer's rise.

\begin{figure}[t]
\centering
\figfile{fig5_long_context.png}
\caption{Loss by token position on PG-19. The dashed line marks the training length.}
\label{fig:long}
\end{figure}

\subsection{Inference efficiency (E5)}
Fig.~\ref{fig:bench} and Table~\ref{tab:bench} give the benchmark on one T4 in fp16. No configuration
ran out of memory up to 32,768 tokens at batch 8; the T4 has 16 GB and the models are small. The
differences are in how the cost grows.

\textbf{Decoding memory.} The Transformer's KV cache holds $2Ld$ fp16 values per layer, 6 GiB at batch
8 and 32K tokens, and its measured peak decoding memory grows from 210.6 MiB at 512 tokens to
6,260.5 MiB at 32K. Mamba's recurrent state is 7.1 MiB at batch 8 whatever the length, and its peak
stays between 114.0 and 116.0 MiB (the weights alone are 89.4 MiB). The \texttt{chunked} backend
gives the same decoding memory and latency as the \texttt{cuda} kernel.

\textbf{Decoding latency.} Mamba with the \texttt{cuda} kernel takes 21.9 to 23.9 ms per new token at
every length and both batch sizes. The Transformer starts faster, 9.5 to 9.9 ms at 512 tokens, and
slows down as the cache grows: 15.1 ms (batch 1) and 24.6 ms (batch 8) at 16K, 28.1 and 46.5 ms at
32K. The two cross between 8K and 16K tokens at batch 8 and between 16K and 32K at batch 1. At a
1,024-token context, the length we trained at, the Transformer decodes 1.7 to 1.8 times faster. We did
not profile why Mamba's step is slower at short contexts; it has twice as many layers to step through.

\textbf{Prefill.} With the \texttt{cuda} kernel, Mamba's prefill throughput at batch 1 stays at 41,000
to 42,000 tokens/s from 1K to 16K tokens and is 31,600 at 32K, while the Transformer's falls from
105,000 tokens/s at 1K to 21,000 at 32K as the $O(L^2)$ attention takes over. The Transformer prefills
faster up to 8K tokens; at 32K and batch 1, Mamba takes 1.04 s and the Transformer 1.53 s. Peak
prefill memory grows with length for both, 910 MiB against 1,319 MiB at 32K and batch 1. Our \texttt{chunked} scan gives the same
outputs but prefills at about 1,550 tokens/s, 20 times slower than the kernel (21.2 s at 32K, batch 1),
and its batch-8 points above 4K tokens were skipped under the 20-second rule. The architecture's
advantage at long context is therefore real with a good kernel and absent without one.

\begin{table}[t]
\centering
\caption{Decoding on one T4, fp16, batch 8: per-token latency and peak memory. Mamba with the
\texttt{cuda} kernel.}
\label{tab:bench}
\begin{tabular}{@{}rrrrr@{}}
\toprule
& \multicolumn{2}{c}{Transformer} & \multicolumn{2}{c}{Mamba (\texttt{cuda})} \\
Context & ms/token & MiB & ms/token & MiB \\
\midrule
512 & 9.9 & 210.6 & 22.5 & 114.0 \\
1,024 & 12.9 & 306.6 & 23.1 & 114.0 \\
4,096 & 12.9 & 882.8 & 23.3 & 114.2 \\
8,192 & 13.3 & 1,651.0 & 23.2 & 114.5 \\
16,384 & 24.6 & 3,187.5 & 23.5 & 115.0 \\
32,768 & 46.5 & 6,260.5 & 23.9 & 116.0 \\
\bottomrule
\end{tabular}
\end{table}

\begin{figure*}[t]
\centering
\figfile[width=\textwidth]{fig6_efficiency.png}
\caption{Prefill latency, peak prefill memory (including the cache), per-token decode latency and peak
decoding memory against length. Line style shows the Mamba scan backend.}
\label{fig:bench}
\end{figure*}

\subsection{What selectivity does (E6)}
Fig.~\ref{fig:horizon} shows, for each layer of the seed-0 Mamba, the median memory horizon
$1/(\Delta_t|A|)$ over the state channels and the 256 analysed tokens. In 22 of the 24 layers the
median horizon is 3 to 31 tokens; the first layer is the shortest (3 tokens, mean $\Delta$ of 0.061),
and the mean $\Delta$ per layer otherwise lies between 0.017 and 0.049. Two layers behave differently:
layer 14 has a median horizon of 80 tokens and layer 18 of 225 tokens, with mean $\Delta$ of 0.0046
and 0.0012. Most of the network therefore summarises the last few dozen tokens, and a small number of
layers carry information across a few hundred, still well under the 1,024-token training window. This
fits E4: if most channels forget within a few hundred tokens, extra context far back cannot help, and
Mamba's loss indeed stops improving past 1K without breaking down. These are medians over 256 tokens
of one model, so we read them as a description, not a proof. The
per-token $\Delta$ heatmap and the Transformer's attention maps are in the repository
(\texttt{results/analysis/S/}); we do not report per-token statistics because the analysed window
contains no end-of-document token, so the question of which tokens trigger a reset cannot be answered
from it.

\begin{figure}[t]
\centering
\figfile{fig7_memory_horizon.png}
\caption{Median memory horizon $1/(\Delta|A|)$ of the state channels in each layer of the seed-0 Mamba,
on 256 test tokens. Layers are numbered from 0.}
\label{fig:horizon}
\end{figure}

\subsection{Instruction fine-tuning (E7)}
Table~\ref{tab:sft} and Fig.~\ref{fig:sft} compare the seed-0 models before and after fine-tuning.
Fine-tuning lowers the test response loss from 2.993 to 2.201 for the Transformer and from 2.975 to
2.163 for Mamba. The gap between the two models is 0.018 before and 0.038 after, in the same direction as
in pretraining (0.046 on the full test split). ROUGE-L on 100 test prompts rises from 0.027 to 0.212
and from 0.000 to 0.222; before fine-tuning, no answer ended with the end-of-document token and every
one ran to the 128-token limit, afterwards 54\% and 55\% stop on their own at 76 to 77 tokens on
average. The pretrained Mamba scores exactly 0 because under greedy decoding it repeats the \texttt{\#}
of the prompt template for all 128 tokens. The pretrained Transformer mostly repeats
``\texttt{\#\#\# Response:}'' instead and scores 0.027.

The sample answers (all in \texttt{results/sft/samples.md}) show what 47M parameters and 340M tokens
buy. Both tuned models answer ``What is the capital of France?'' with ``The capital of France is
Paris.'' Both start their photosynthesis answer with a correct first sentence and then add a muddled
second one. Both fail at arithmetic: to
``What is 12 plus 7?'' Mamba answers ``The 12 plus 7 is equal to 12.'' and the Transformer repeats
``12 + 7 = 12 + 7 = \dots'' until cut off. Asked for three ways to save water, both loop on brooms and
hoses (``use a broom instead of a hose to clean the air''). Asked to summarise a short text about the
sun, the Transformer copies the first sentence of the text and Mamba writes ``plants use it to grow and
grow''. We see no consistent difference in answer quality between the two architectures at this
scale; both learned the answer format well and the facts poorly.

\begin{table}[t]
\centering
\caption{Instruction fine-tuning, seed-0 models, test split (1,075 examples; ROUGE-L and stopping rate
on the first 100 prompts, greedy decoding, at most 128 new tokens).}
\label{tab:sft}
\begin{tabular}{@{}llrrrr@{}}
\toprule
Model & Stage & Resp. loss & Ppl & ROUGE-L & Stops \\
\midrule
Transformer & pretrained & 2.993 & 19.95 & 0.027 & 0\% \\
Transformer & tuned & 2.201 & 9.03 & 0.212 & 54\% \\
Mamba & pretrained & 2.975 & 19.60 & 0.000 & 0\% \\
Mamba & tuned & \textbf{2.163} & \textbf{8.70} & \textbf{0.222} & 55\% \\
\bottomrule
\end{tabular}
\end{table}

\begin{figure}[t]
\centering
\figfile{fig8_sft.png}
\caption{Test response loss and ROUGE-L before and after instruction fine-tuning.}
\label{fig:sft}
\end{figure}

% ---------------------------------------------------------------------------------------------------
\section{Discussion}
\textbf{The research question.} At 47M parameters, the same data and 340M tokens, the Mamba-style model
reaches a lower test loss than the Transformer (3.484 against 3.531), and the difference is many times
the seed spread. The Mamba paper~\cite{gu2023mamba} reports the same ordering from 130M parameters up;
we find it holds at a third of that size and about 7 tokens per parameter. In inference, Mamba's
decoding memory is flat in the context length and its per-token latency is flat, where the
Transformer's both grow, but at the 1K context we trained at the Transformer decodes faster. Which
model is ``more efficient'' depends on the context length one cares about.

\textbf{Where the loss comes from.} The recall split explains the shape of the result. The Transformer
is 0.37 nats better on the 16.8\% of tokens whose answer sits earlier in the window, which is the
behaviour Zoology~\cite{arora2023zoology} predicts for attention. Mamba is 0.13 nats better on the
other 83\%, and because those are five times as many, it wins overall. The two advantages are of
different kinds: attention looks a value up exactly; the recurrence compresses the whole past into a
state of 16 numbers per channel and apparently makes better use of it for ordinary next-token
prediction at this size. MQAR does not separate the two: both solve it at widths 128 and 256, each at
one learning rate, and neither at width 64.

\textbf{What selectivity buys.} Mamba-LTI has the same size and training and the same recurrence,
minus the dependence of $\Delta$, $B$ and $C$ on the input. It loses 0.65 nats on recall-hit tokens
and only 0.08 elsewhere, needs a higher learning rate and two to three times more steps on selective
copying,
and never solves MQAR. So the input-dependent step is what lets the recurrence decide what to store
and when to overwrite it; without it the model is a reasonable language model for local prediction and
a poor one for copying. The memory horizons support this: most layers forget within a few dozen tokens,
and two layers hold information for 80 and 225 tokens.

\textbf{Architecture against implementation.} Two of our efficiency results depend on the kernel and
not on the model. With the \texttt{cuda} kernel, Mamba prefills 32K tokens faster than the Transformer;
with our \texttt{chunked} scan it is 14 times slower than the Transformer at the same length. And on
the T4, Mamba's training runs took twice the Transformer's time for 9\% fewer FLOPs, so at equal GPU
hours the Transformer gave the lower loss. The architecture's asymptotic advantages are real, but at
47M parameters and on a 2018 GPU the Transformer's mature kernels decide the practical outcome.

\textbf{Fine-tuning.} Instruction tuning does not change the comparison. Both models drop their
response loss by about 0.8 nats, both learn the answer format, and the 0.04-nat gap after tuning has
the same sign and about the same size as before. Their mistakes are also the same kind: repetition and
wrong facts.

\section{Limitations}
\begin{itemize}
  \item \textbf{Scale and budget.} The models have 47M parameters and see 340M tokens, about 7 per
        parameter, below the compute-optimal ratio of about 20. Published Mamba results start at 130M
        parameters with much larger budgets, so our trends may not hold at larger scale. A pair of about
        98M parameters was configured but not trained; the smoke test estimated about 104 to 105 hours
        per run.
  \item \textbf{Seeds.} Three seeds for the Transformer and Mamba and two for Mamba-LTI. Differences
        smaller than the seed spread should not be called significant.
  \item \textbf{Learning rate.} Tuned with 30M-token runs; Mamba's lowest sweep loss was at the top of
        the grid, and Mamba-LTI was not tuned separately.
  \item \textbf{One corpus, one tokenizer.} Perplexities compare only models that share this tokenizer
        and test set.
  \item \textbf{Free hardware.} The T4 has no native bf16, so training used fp16. Efficiency and
        training-time results describe these implementations on this GPU. The equal-time comparison
        reads Mamba's loss from the middle of a longer run, which understates what a run planned for
        four hours would reach.
  \item \textbf{Synthetic tasks.} One seed, one sequence length for MQAR, and an outcome that flips
        with the learning rate. The 4,000-step runs show how much the conclusion depends on training
        length.
  \item \textbf{Proxies.} The recall-hit split counts repeated token pairs, which is only one form of
        recall. PG-19 books differ from the web training text, which affects both models.
  \item \textbf{Instruction fine-tuning.} One seed and one dataset, whose answers were written by
        another model. ROUGE-L measures word overlap with one reference answer, not correctness, and
        answers are cut at 128 new tokens.
\end{itemize}

\section{Conclusion}
We trained a Transformer, a Mamba and a non-selective Mamba of about 47M parameters on the same 340M
tokens of FineWeb-Edu, on free GPUs, and compared them on quality, recall, long context, inference cost
and instruction following. Mamba reached the lower test loss (3.484 against 3.531) by being better on
the 83\% of tokens that cannot be copied from the context, while the Transformer stayed better on the
17\% that can. Removing selectivity cost 0.18 nats, and per token the damage was eight times larger
on the copyable tokens (0.65 against 0.08). Mamba's
decoding memory did not grow with context and its loss degraded far less past the training length,
but the Transformer trained twice as fast on a T4, and without the official kernel Mamba's scan was
20 times slower. Instruction tuning improved both models by the same amount and left the ordering
unchanged.

\section*{Code and Data}
Code, configurations, the Kaggle notebook and the per-run logs:
\url{https://github.com/Gagan-1718/afml-transformer-vs-state-space-models}. FineWeb-Edu revision
\texttt{87f09149ef4734204d70ed1d046ddc9ca3f2b8f9}; the tokenizer's SHA-256 and the split sizes are in
\texttt{results/data/meta.json}.

\begin{thebibliography}{99}

\bibitem{vaswani2017attention}
A.~Vaswani, N.~Shazeer, N.~Parmar, J.~Uszkoreit, L.~Jones, A.~N. Gomez, {\L}.~Kaiser, and
I.~Polosukhin, ``Attention is all you need,'' in \emph{Advances in Neural Information Processing
Systems (NeurIPS)}, 2017.

\bibitem{gu2022s4}
A.~Gu, K.~Goel, and C.~R{\'e}, ``Efficiently modeling long sequences with structured state spaces,''
in \emph{International Conference on Learning Representations (ICLR)}, 2022.

\bibitem{gu2023mamba}
A.~Gu and T.~Dao, ``Mamba: Linear-time sequence modeling with selective state spaces,''
arXiv:2312.00752, 2023.

\bibitem{arora2023zoology}
S.~Arora, S.~Eyuboglu, A.~Timalsina, I.~Johnson, M.~Poli, J.~Zou, A.~Rudra, and C.~R{\'e}, ``Zoology:
Measuring and improving recall in efficient language models,'' arXiv:2312.04927, 2023.

\bibitem{dao2024mamba2}
T.~Dao and A.~Gu, ``Transformers are SSMs: Generalized models and efficient algorithms through
structured state space duality,'' in \emph{International Conference on Machine Learning (ICML)}, 2024.

\bibitem{zhang2019rmsnorm}
B.~Zhang and R.~Sennrich, ``Root mean square layer normalization,'' in \emph{Advances in Neural
Information Processing Systems (NeurIPS)}, 2019.

\bibitem{su2021roformer}
J.~Su, Y.~Lu, S.~Pan, A.~Murtadha, B.~Wen, and Y.~Liu, ``RoFormer: Enhanced transformer with rotary
position embedding,'' arXiv:2104.09864, 2021.

\bibitem{shazeer2020glu}
N.~Shazeer, ``GLU variants improve transformer,'' arXiv:2002.05202, 2020.

\bibitem{penedo2024fineweb}
G.~Penedo, H.~Kydl{\'i}{\v{c}}ek, L.~Ben~Allal, A.~Lozhkov, M.~Mitchell, C.~Raffel, L.~Von~Werra, and
T.~Wolf, ``The FineWeb datasets: Decanting the web for the finest text data at scale,''
arXiv:2406.17557, 2024.

\bibitem{rae2020compressive}
J.~W. Rae, A.~Potapenko, S.~M. Jayakumar, C.~Hillier, and T.~P. Lillicrap, ``Compressive transformers
for long-range sequence modelling,'' in \emph{International Conference on Learning Representations
(ICLR)}, 2020.

\bibitem{loshchilov2019adamw}
I.~Loshchilov and F.~Hutter, ``Decoupled weight decay regularization,'' in \emph{International
Conference on Learning Representations (ICLR)}, 2019.

\bibitem{hoffmann2022chinchilla}
J.~Hoffmann \emph{et~al.}, ``Training compute-optimal large language models,'' arXiv:2203.15556, 2022.

\bibitem{wei2022flan}
J.~Wei, M.~Bosma, V.~Y. Zhao, K.~Guu, A.~W. Yu, B.~Lester, N.~Du, A.~M. Dai, and Q.~V. Le, ``Finetuned
language models are zero-shot learners,'' in \emph{International Conference on Learning Representations
(ICLR)}, 2022.

\bibitem{taori2023alpaca}
R.~Taori, I.~Gulrajani, T.~Zhang, Y.~Dubois, X.~Li, C.~Guestrin, P.~Liang, and T.~B. Hashimoto,
``Stanford Alpaca: An instruction-following LLaMA model,''
\url{https://github.com/tatsu-lab/stanford_alpaca}, 2023.

\bibitem{wang2023selfinstruct}
Y.~Wang, Y.~Kordi, S.~Mishra, A.~Liu, N.~A. Smith, D.~Khashabi, and H.~Hajishirzi, ``Self-Instruct:
Aligning language models with self-generated instructions,'' in \emph{Annual Meeting of the Association
for Computational Linguistics (ACL)}, 2023.

\bibitem{lin2004rouge}
C.-Y. Lin, ``ROUGE: A package for automatic evaluation of summaries,'' in \emph{Text Summarization
Branches Out (ACL Workshop)}, 2004.

\end{thebibliography}

\end{document}
